\section{Numerical experiments}


In this section, we analyze the convergence properties of the finite element approximation to the proposed problem. We consider a sequence of finite element conformal and shape-regular partitions of the computational domain $\mathcal{T}_h$ with characteristic mesh size $h$. 
Let us define the discrete spaces associated to $\mathcal{T}_h$
$$
\CG_{k}(\mathcal{T}_h) = \{v \in C^0(\Omega), v|_K \in \mathbb{P}_k(K)~\forall K \in \mathcal{T}_h\}
$$
and 
$$
\DG_{k}(\mathcal{T}_h) = \{v, v|_K \in \mathbb{P}_k(K)~\forall K \in \mathcal{T}_h\}
$$
where $\mathbb{P}_k(K)$ is the space of polynomials of degree $k$ on $K$.

Two different choices of finite element spaces are then considered to approximate
the optimization problem, namely, a \textbf{natural formulation}, in which we take 
\begin{equation}\label{eq:natural_func_spaces}
	\left(\uP_{h}, \eP_{h}, \sP_{h}, \lP_{h}, \mP_{h}\right) \in \CG_{k+1} \times \left[\DG_{k}\right]^{d_{\Omega}} \times \left[\DG_{k}\right]^{d_{\Omega}} \times \CG_{k+1} \times \left[\DG_{k}\right]^{d_{\Omega}} 
\end{equation} 
and an \textbf{equal-order formulation}, in which we take
\begin{equation}\label{eq:equal_order_func_spaces}
\left(\uP_{h}, \eP_{h}, \sP_{h}, \lP_{h}, \mP_{h}\right)  \in \CG_{k+1} \times \left[\CG_{k+1}\right]^{d_{\Omega}} \times \left[\CG_{k+1}\right]^{d_{\Omega}} \times \CG_{k+1} \times \left[\CG_{k+1}\right]^{d_{\Omega}}
\end{equation}
In the numerical tests we take $k \in \{0, 1, 2\}$ for both formulations, corresponding to the use of linear, quadratic or cubic polynomials for the scalar fields $\uP$ and $\lP$ in both formulations. Indeed, the Natural formulation
considers one degree lower for the vector fields $\sP$, $\eP$ and $\mP$, whereas for the equal-order formulation, the degree is the same for all fields. 
To assess the convergence properties of the method we consider both convex and non-convex geometries in two spatial dimensions. The convergence behavior is assessed by computing the error with respect to manufactured solutions.
Finally, for the equal-order formulation, we consider three variants: a Unstabilized method, the Fully Stabilized method, and the Minimally Stabilized method.
In the Unstabilized method, all stabilization parameters are set to zero.
In the Fully Stabilized method, the parameters are chosen according to equation (\ref{eq:paramsopt}). For the Minimally Stabilized method, all parameters are set to zero except for $\alpha$ and $\eta$, which represent the minimum requirements for stability, as indicated in equation (\ref{eq:Bhparams}).

\subsection{Test case 1: a convex domain}\label{sec:test_case1}

Let us first describe the procedure for constructing the manufactured solution for this test problem. First, consider the smooth field $\uP$ defined by:
\begin{equation}
	\uP:\Omega \to \mathbb{R},\quad \uP(x,y) = \cos(\pi x)\cos(\pi y),
\end{equation}
whose gradient is given by:
\begin{equation*}
	\begin{aligned}
	\eP: \Omega \to \mathbb{R}^{2}, \quad \eP(x,y) = -\pi \, [\sin(\pi x)\cos(\pi y), \cos(\pi x)\sin(\pi y)]^{\top}. \\
	\end{aligned}
\end{equation*}
Also define 
$$
\lP: \Omega \rightarrow \mathbb{R},\qquad \lP(x,y) = -\frac{1}{40\,\pi}\,\left[\cos\left(2\,\pi\,x \right) + \cos\left(2\,\pi\,y \right)\right] \\ \\
$$
and
\begin{equation}
\sP:\Omega \rightarrow \mathbb{R}^2, \qquad \sP = -\dfrac{\eP}{\|\eP\|}\, \left(\|\eP\| - \dfrac{1}{40}\,\|\eP\|^{3}\right) \label{eq:constlaw}
\end{equation}
At this stage, solely for the purpose of assessing the convergence properties of the discrete formulations under general conditions, it is convenient to introduce an additional source term on the right-hand side of eq. (\ref{eq:dif1}), i.e.,
\begin{equation}
	\zeta \lP  + \nabla \cdot \mP  =  f \label{eq:balance_with_f}
\end{equation}
This amounts to adding the term $-({f , \uP})$ to the Lagrangian (\ref{eq:lagrangian}).
Notice that the addition of this linear term does not affect the convergence analysis from previous sections.
By examining the optimality conditions in their strong form (\ref{eq:dif1})-(\ref{eq:dif5}),
we proceed as follows to construct the manufactured solution for the remaining fields.
First take
$$
\etP:\Omega \rightarrow \mathbb{R}^2, \qquad \etP(x,y) = \eP(x,y) + \frac{1}{20} \left[\sin\left(4\,\pi\,x \right), \sin\left(4\,\pi\,y\right) \right]^{\top}
$$
as being a perturbation of $\eP$, from which $\mP$ can be obtained using eq. (\ref{eq:dif2}).
Now, $\stP$ is computed using eq. (\ref{eq:dif3}) and finally, the source 
terms $f$ and $q$ are obtained after substituting the proposed fields 
into eqs. (\ref{eq:dif1}) and (\ref{eq:dif4}). 
In the numerical experiments we take $\kappa = 1$ and $\zeta = 1$.
Notice that, in this way, the fields $\eP$ and $\sP$ satisfy a type of 
constitutive relationship, which according to (\ref{eq:constlaw}) is essentially
a linear law plus a cubic perturbation term, whereas the values taken by the fields 
$\etP$ and $\stP$ at different points throughout the domain do not necessarily do so. 
This mimics practical scenarios in which these are obtained from experimental 
measurements or synthetic data. This is illustrated in part (a) of figure \ref{fig:all_fields_ref} which shows
a plot of the norm of $\sP$ as a function of the norm of $\eP$, along with their counteparts $\stP$ and $\etP$ evaluated over a computational mesh. 
Notice that the exact flux-gradient relationship is depicted as a strictly monotonic curve (thick black line), whereas the fields $\stP$ and $\etP$ (blue dots) exhibit a certain dispersion around the former. Contours of the different fields are also shown in the
figure considering the computational domain $\Omega$ to be the unit square $[0,1]^2$. The boundary of $\Omega$ is defined as $\partial \Omega = \Gamma_{\text{left}} \cup \Gamma_{\text{right}} \cup \Gamma_{\text{bottom}} \cup \Gamma_{\text{top}}$.
The problem is solved imposing the following boundary conditions
\begin{equation}
	\left\lbrace
	\begin{array}{rclll}
		\sP \cdot \mathbf{\check{n}} 			&=& \mathsf{g}_{\sP} 			&\mbox{on}&~\Gamma_{\text{bottom}} \cup \Gamma_{\text{top}} \\
		\mP \cdot \mathbf{\check{n}} &=& \mathsf{g}_{\mP} &\mbox{on}&~\Gamma_{\text{bottom}} \cup \Gamma_{\text{top}} \\
		\uP &=& \mathsf{g}_{\uP} &\mbox{on}&~\Gamma_{\text{left}} \cup \Gamma_{\text{right}} \\
		\lP &=& \mathsf{g}_{\lP} &\mbox{on}&~\Gamma_{\text{left}} \cup \Gamma_{\text{right}} \\
	\end{array}
	\right.
\end{equation}
where the values in the right hand side are computed from the manufactured fields.
\begin{figure}[h!]
	\centering
	\caption{Figure (a) presents the functional relationship between the norm of the flux and the gradient $\|\sP\|~\mbox{vs.}~\|\eP\|$, and $\|\stP\|~\mbox{vs.}~\|\etP\|$, respectively. Figure (b) to (f) shows
		$\uP$, $\lP$, $\eP$, $\sP$, and $\mP$.}
	%\fontsize{11pt}{10pt}\selectfont 
	\def\svgwidth{1.0\columnwidth}
	\input{Figures/PDFs/all_fields_ref.pdf_tex}
	\label{fig:all_fields_ref}
\end{figure}

Moving to the results, Figure \ref{fig:convergence_scalar_linear} shows the $L^{2}$- and $H^{1}$-error norms for the scalar fields $\uP$ and $\lP$ with $k=0$, for both the Natural formulation (subfigures (a) and (b)) and the equal-order formulations (subfigures (c)-(f)). In the latter case, the unstabilized, minimally stabilized, and fully stabilized variants are considered. In all cases, the convergence rates are consistent with the expected theoretical error estimate, which is of order $\mathcal{O}(h)$ in the $H^{1}$-norm (see Theorem~2). Also, the error in the $ L^{2} $-norm converges with order $ \mathcal{O}(h^2) $.
No significative differences are observed between the different equal order variants.
Figure~\ref{fig:convergence_vector_linear} shows the corresponding errors for the vector fields
$\sP$, $\eP$ and $\mP$ in the $ L^{2} $-norm. For the Natural formulation linear convergence is observed for all fields, whereas the equal order formulation exhibits superlinear behavior, almost reaching $ \mathcal{O}(h^2) $. 

The corresponding numerical results for $\text{P}_2$-elements ($k=1$) and $\text{P}_3$-elements ($k=2$) are shown in Figures~\ref{fig:convergence_scalar_quadratic} and ~\ref{fig:convergence_scalar_cubic}
for the scalar fields and Figures~\ref{fig:convergence_vector_quadratic} and ~\ref{fig:convergence_vector_cubic}
for the vector fields. Notice the superior behavior of the Minimally and Fully stabilized 
variants with respect to the Unstabilized case, reaching, respectively, second and third order convergence rates for the scalar fields $\uP$ and $\lP$ and for the vector
fields $\eP$ and $\mP$. The convergence rate is better than expected
for the flux $\sP$ attaining $ \mathcal{O}(h^3) $ for quadratics and $ \mathcal{O}(h^4) $ 
for cubics. Notably, in all cases, the Fully stabilized case exhibited the lowest error values.

Although we do not have theoretical error estimates in the $ \text{H(div)} $-norm, 
convergence rates are also reported. Figures \ref{fig:convergence_vector_linear_hdiv}, 
\ref{fig:convergence_vector_quadratic_hdiv} and \ref{fig:convergence_vector_cubic_hdiv}
show the results for $k=0$, $k=1$ and $k=2$, respectively.
Note that the Fully stabilized formulation yielded the most favorable results in all three cases. For $ \text{P}_1 $-elements, no significant differences were observed in the convergence behavior of the flux $ \sP_{h} $ across the Unstabilized, Minimally stabilized, and Fully stabilized formulations, with the corresponding error curves largely overlapping. In contrast, for $ \eP $ and $ \mP $, the Fully stabilized formulation exhibited sub-linear convergence, unlike the Unstabilized case. However, it achieved lower error magnitudes. For $ \text{P}_2 $-elements, the Fully stabilized formulation exhibited the best performance. In particular, it was the only one to attain second-order convergence for the flux $ \sP $, and exhibited the smallest errors for $ \eP_{h} $ and $ \mP_{h} $. A similar trend was observed for $ \text{P}_3 $-elements. 

It is worth mentioning that additional numerical experiments were carried out to evaluate the effect of each individual stabilization parameter in isolation and compared to the Unstabilized, Minimally stabilized, and Fully stabilized cases. The results are not reported here for the sake of brevity, however, they displayed only minor variations in magnitude and exhibited convergence rates consistent with those of the main reference cases. Furthermore, numerical experiments confirmed that the convergence behavior is rather independent of the reaction coefficient $ \zeta $, as similar results were obtained when varying $\zeta$ over several orders of magnitude.

In what follows, we illustrate the stability properties of the approximation of the Natural formulation and the variants of the equal-order formulation by evaluating the difference \( \uP_{h} - \uP \) on a sufficiently fine mesh. Figure~\ref{fig:checkerboard_pattern} presents the corresponding results. We observe that no spurious modes can be identified in either the Natural formulation or the Fully stabilized variant. In the Minimally stabilized variant, a subtle checkerboard pattern appears near the boundaries. In contrast, in the Unstabilized case, spurious modes are manifested throughout the entire domain. This behavior is explained by the fact that the space \( N_{h} \subset Y_{h} := L_{h} \times M_{h} \), defined as
$$
N_h = \left\{ y_{h} \in Y_{h} \,\middle|\,
b(x_h, y_{h}) = 0 \quad \forall x_h \in X_h \right\},
$$
is nontrivial.
%Checkerboard pattern
\begin{figure}[h!]
	\centering
	\caption{Figures~(a)--(d) show the difference $\uP_{h} - \uP$ for the Natural formulation, and for the equal-order Unstabilized, Fully stabilized, and Minimally stabilized formulations, respectively. }
	%\fontsize{12pt}{10pt}\selectfont 
	\def\svgwidth{1.0\columnwidth}
	\input{Figures/PDFs/checkerboard_pattern.pdf_tex}
	\label{fig:checkerboard_pattern}
\end{figure}
As a final comment, we emphasize that the manufactured fields $ \eP $ and $ \sP $ were constructed to be physically admissible in the sense that they do not violate the second law of thermodynamics -- the \textit{Clausius–Duhem inequality} \cite{Reddy2013,Gurtin2013} --. Accordingly, the following condition must hold:
$$
-\sP(\mathbf{x}) \cdot \eP(\mathbf{x}) \geq 0, \quad \forall \mathbf{x} \in \Omega
$$ 
However, even though the discrete approximations $ \eP_{h} $ and $ \sP_{h} $ converge to the manufactured fields in norm, this does not guarantee that they satisfy the second law pointwise. Therefore, for all the problems addressed in this work, we have explicitly verified the physical admissibility between flux and gradient fields, and no violations were identified. 


%where the bilinear form $b(x, y)$ is:
%\begin{equation}
%	b(x, y) = (\mu, e - \nabla u) - (\nabla \lambda, s) + (\lambda, \zeta u).
%\end{equation}
%
%Thus:
%\begin{equation}
%	N_h = \left\{ (\lambda_h, \mu_h) \in L_h \times M_h ~\middle|~
%	\begin{aligned}
%		&(\mu_h, e_h - \nabla u_h) - (\nabla \lambda_h, s_h) + (\lambda_h, \zeta u_h) = 0 \\
%		&\quad \forall (u_h, e_h, s_h) \in U_h \times E_h \times S_h
%	\end{aligned}
%	\right\}.
%\end{equation}

 


%=============================================
% k = 0

\begin{figure}[h!]
	\centering
	\caption{Convergence of scalar fields $\uP$ and $\lP$ with $k=0$. Figures~(a)--(b) show results for $\uP$, and (c)--(d) for $\lP$ in both Natural and equal-order formulations, evaluated in the $L^{2}$ and $H^{1}$ norms.   Three stabilization configurations are compared: Unstabilized, Fully stabilized, and Minimally stabilized ($\alpha$, $\eta$).
	}
	%\fontsize{11pt}{10pt}\selectfont 
	\def\svgwidth{1.0\columnwidth}
	\input{Figures/PDFs/convergence_scalar_linear_new.pdf_tex}
	\label{fig:convergence_scalar_linear}
\end{figure}


\begin{figure}[h!]
	\centering
	\caption{Convergence of vector fields $\sP$, $\eP$, and $\mP$ with $k=0$, evaluated in the $L^{2}$-norm for both Natural and equal-order formulations. Three stabilization configurations are compared: Unstabilized, Fully stabilized, and Minimally stabilized ($\alpha$, $\eta$).
	}
	%\fontsize{11pt}{10pt}\selectfont 
	\def\svgwidth{1.0\columnwidth}
	\input{Figures/PDFs/convergence_vector_linear_new.pdf_tex}
	\label{fig:convergence_vector_linear}
\end{figure}


\begin{figure}[h!]
	\centering
	\caption{Convergence of vector fields $\sP$, $\eP$, and $\mP$ with $k=0$ in the equal-order formulation, evaluated in the $\text{Hdiv}$-norm. Figures~(a)--(c) present results for three stabilization configurations: Unstabilized, Fully stabilized, and Minimally stabilized ($\alpha$, $\eta$).
	}
	%\fontsize{11pt}{10pt}\selectfont 
	\def\svgwidth{1.0\columnwidth}
	\input{Figures/PDFs/convergence_vector_linear_hdiv.pdf_tex}
	\label{fig:convergence_vector_linear_hdiv}
\end{figure}



%=============================================
% k = 1

\begin{figure}[h!]
	\centering
	\caption{Convergence of scalar fields $\uP$ and $\lP$ with $k=1$. Figures~(a)--(b) show results for $\uP$, and (c)--(d) for $\lP$ in both Natural and equal-order formulations, evaluated in the $L^{2}$ and $H^{1}$ norms.   Three stabilization configurations are compared: Unstabilized, Fully stabilized, and Minimally stabilized ($\alpha$, $\eta$).}
	%\fontsize{11pt}{10pt}\selectfont 
	\def\svgwidth{1.0\columnwidth}
	\input{Figures/PDFs/convergence_scalar_quadratic_new.pdf_tex}
	\label{fig:convergence_scalar_quadratic}
\end{figure}


\begin{figure}[h!]
	\centering
	\caption{Convergence of vector fields $\sP$, $\eP$, and $\mP$ with $k=1$, evaluated in the $L^{2}$-norm for both Natural and equal-order formulations. Three stabilization configurations are compared: Unstabilized, Fully stabilized, and Minimally stabilized ($\alpha$, $\eta$).}
	%\fontsize{11pt}{10pt}\selectfont 
	\def\svgwidth{1.0\columnwidth}
	\input{Figures/PDFs/convergence_vector_quadratic_new.pdf_tex}
	\label{fig:convergence_vector_quadratic}
\end{figure}

\begin{figure}[h!]
	\centering
	\caption{Convergence of vector fields $\sP$, $\eP$, and $\mP$ with $k=1$ in the equal-order formulation, evaluated in the $\text{Hdiv}$-norm. Figures~(a)--(c) present results for three stabilization configurations: Unstabilized, Fully stabilized, and Minimally stabilized ($\alpha$, $\eta$).}
	%\fontsize{11pt}{10pt}\selectfont 
	\def\svgwidth{1.0\columnwidth}
	\input{Figures/PDFs/convergence_vector_quadratic_hdiv.pdf_tex}
	\label{fig:convergence_vector_quadratic_hdiv}
\end{figure}




%=============================================
% k = 2
\begin{figure}[h!]
	\centering
	\caption{Convergence of scalar fields $\uP$ and $\lP$ with $k=2$. Figures~(a)--(b) show results for $\uP$, and (c)--(d) for $\lP$ in both Natural and equal-order formulations, evaluated in the $L^{2}$ and $H^{1}$ norms.   Three stabilization configurations are compared: Unstabilized, Fully stabilized, and Minimally stabilized ($\alpha$, $\eta$).}
	%\fontsize{11pt}{10pt}\selectfont 
	\def\svgwidth{1.0\columnwidth}
	\input{Figures/PDFs/convergence_scalar_cubic_new.pdf_tex}
	\label{fig:convergence_scalar_cubic}
\end{figure}



\begin{figure}[h!]
	\centering
	\caption{Convergence of vector fields $\sP$, $\eP$, and $\mP$ with $k=2$, evaluated in the $L^{2}$-norm for both Natural and equal-order formulations. Three stabilization configurations are compared: Unstabilized, Fully stabilized, and Minimally stabilized ($\alpha$, $\eta$).}
	%\fontsize{11pt}{10pt}\selectfont 
	\def\svgwidth{1.0\columnwidth}
	\input{Figures/PDFs/convergence_vector_cubic_new.pdf_tex}
	\label{fig:convergence_vector_cubic}
\end{figure}

\begin{figure}[h!]
	\centering
	\caption{Convergence of vector fields $\sP$, $\eP$, and $\mP$ with $k=2$ in the equal-order formulation, evaluated in the $\text{Hdiv}$-norm. Figures~(a)--(c) present results for three stabilization configurations: Unstabilized, Fully stabilized, and Minimally stabilized ($\alpha$, $\eta$).}
	%\fontsize{11pt}{10pt}\selectfont 
	\def\svgwidth{1.0\columnwidth}
	\input{Figures/PDFs/convergence_vector_cubic_hdiv.pdf_tex}
	\label{fig:convergence_vector_cubic_hdiv}
\end{figure}


\cleardoublepage
\newpage
\subsection{Test case 2: a non-convex domain}\label{sec:test_case2}

In this case we consider the classical problem with a re-entrant corner, whose geometry is depicted in 
Figure \ref{fig:geometries_non_convex} for three different external angles, namely, $\phi =\pi/4,~\pi/2$ and $3\pi/4$. The corner is located at the origin $(0, 0)$. 
We manufacture a solution by taking the scalar field $u$ as being:
\begin{equation*}
	\begin{aligned}
	&\uP: \Omega \rightarrow \mathbb{R}, \quad \uP(x,y) = r^{\nu}\,\sin(\nu\,\theta) \\
	&\mbox{where}~r = \sqrt{\left(x^{2}+y^{2}\right)}, \quad \theta = \arctan{\left(\frac{y}{x}\right)}, \quad  \nu = \frac{\pi}{\psi},~\mbox{with}~\psi = 2\pi-\phi
	\end{aligned}
\end{equation*}
The corresponding gradient field is then given by
\begin{equation*}
	\begin{aligned}
	\eP: \Omega \rightarrow \mathbb{R}^{2}, \quad \eP(x,y) = \nu\,r^{\nu-2}\left[x\,\sin(\nu\,\theta)-y\,\cos(\nu\,\theta),y\,\sin(\nu\,\theta)+x\,\cos(\nu\,\theta)\right]^{\top}
	\end{aligned}
\end{equation*}
In this setting, the expressions for $\lP$ and $\etP$ are taken to be the same as those used in the preceding numerical experiment. For $\sP$ we simply consider the linear relation $\sP = -\eP$.
Also, we assume pure Dirichlet boundary conditions for $\uP$ and $\lP$ based on the manufactured solution.
\begin{figure}[h!]
	\centering
	\caption{Illustration of the three non-convex geometries considered.}
	\fontsize{11pt}{10pt}\selectfont 
	\def\svgwidth{1.0\columnwidth}
	\input{Figures/PDFs/geometries_non_convex.pdf_tex}
	\label{fig:geometries_non_convex}
\end{figure}
Figure~\ref{fig:all_fields_ref_non_convex} illustrates the manufactured fields. Again, we also show the behavior of the norm of $\sP$ as a function of the norm of  $\eP$ and their counterparts $ \stP$ and $\etP $.
It is worth noticing that the cubic term considered in the previous experiment has been suppressed so as to not exacerbate the singular behavior of the solution.
\begin{figure}[h!]
	\centering
	\caption{Figure (a) presents the functional relationship between the norm of the flux and the gradient $\|\sP\|~\mbox{vs.}~\|\eP\|$, and $\|\stP\|~\mbox{vs.}~\|\etP\|$, respectively. Figure (b) to (f) show $\uP$, $\lP$, $\eP$, $\sP$, and $\mP$.}
	%\fontsize{11pt}{10pt}\selectfont 
	\def\svgwidth{1.0\columnwidth}
	\input{Figures/PDFs/all_fields_ref_non_convex.pdf_tex}
	\label{fig:all_fields_ref_non_convex}
\end{figure}
For all these geometries, it can be shown that $\uP$ possesses fewer than $\nu+1$ derivatives in $L^{2}$, implying that $\uP \in H^{1}(\Omega)$, but $\uP \notin H^{2}(\Omega)$. Consequently, no improvement in the convergence rate is expected from polynomial refinement and therefore we only use $\text{P}_1$-elements.
In order to better capture the solution, a local mesh refinement near the singularity is employed, as illustrated in Figure~\ref{fig:geometries_non_convex}(a).
We assess the convergence behavior with respect to different corner angles. Specifically, for the values $ \phi \in \left\{ \frac{\pi}{4}, \frac{\pi}{2}, \frac{3\pi}{4} \right\} $, the corresponding values of $ \psi $ are $ \left\{ \frac{7\pi}{4}, \frac{3\pi}{2}, \frac{5\pi}{4} \right\} $. Consequently, the scalar field $ \uP $ belong to the spaces $ H^{11/7 - \varepsilon}(\Omega) $, $ H^{5/3 - \varepsilon}(\Omega) $, and $ H^{9/5 - \varepsilon}(\Omega) $, for all $ \varepsilon > 0 $. Error estimates for finite element schemes in domains with corners have been extensively studied in the literature, for instance by \cite{Babuska1972, Babuska1970}. As a result, the convergence rate of $ \uP $ in both $ L^{2} $-norm  and $ H^{1} $-norms are expected to be restricted by the singular behavior at the reentrant corner and one may anticipate convergence rates of $\mathcal{O}(h^{1+4/7-\varepsilon})$, $ \mathcal{O}(h^{1+2/3-\varepsilon}) $ and $ \mathcal{O}(h^{1+4/5-\varepsilon}) $ for all $ \varepsilon > 0 $ in $L^{2}$-norm, and $\mathcal{O}(h^{4/7-\varepsilon})$, $ \mathcal{O}(h^{2/3-\varepsilon}) $ and $ \mathcal{O}(h^{4/5-\varepsilon}) $ for all $ \varepsilon > 0 $ in $H^{1}$-norm. 
Table~\ref{tab:scalar_errors} summarizes the convergence rates for the scalar fields in both $ L^{2} $- and $ H^{1} $-norms, considering the Natural formulation, the equal-order formulation, and its stabilized variants, across the three corner angles. Similarly, Table~\ref{tab:vector_errors} reports the convergence rates of the vector fields in the $ L^{2} $-norm.
Focusing on the scalar field $ \uP $ in the Natural formulation, we observe that the convergence in the $ H^1 $-norm closely approaches the theoretical upper bounds. In the $ L^2 $-norm, the convergence rate approximately doubles that of the $ H^1 $-norm, as expected from standard interpolation theory. The field $ \lP $ exhibits slightly better convergence than $ \uP $, particularly in the $ H^1 $-norm.
\begin{table}[h!]
	\centering
	\caption{Estimated convergence rates for the scalar fields $ \uP_h $ and $ \lP_h $ in the $ L^2 $- and $ H^1 $-norms using different stabilization schemes and reentrant corner angles $ \psi \in \left\{ \frac{7\pi}{4}, \frac{3\pi}{2}, \frac{5\pi}{4} \right\} $.}
	%\rowcolors{3}{black!10!white}{white}
	\begin{tabular}{l l l l l}
		\toprule
		%\rowcolor{black!30!white} 
		Formulation & Norms & $\psi = \frac{7\pi}{4}$ & $\psi = \frac{3\pi}{2}$ & $\psi = \frac{5\pi}{4}$ \\ [2pt]
		\hline
		\multirow{4}{*}{Natural}				 & $\|u-\uP_{h}\|_{L^{2}}$ & 1.14  & 1.44  & 1.65  \\
		                                 & $\|u-\uP_{h}\|_{H^{1}}$ & 0.57 & 0.72  & 0.81 \\
		                                 & $\|\lP-\lP_{h}\|_{L^{2}}$ & 1.22  & 1.68 & 1.97  \\
		                                 & $\|\lP-\lP_{h}\|_{H^{1}}$ & 0.99  & 0.99 & 0.99  \\
		\hline
		\multirow{4}{*}{Unstabilized}    & $\|u-\uP_{h}\|_{L^{2}}$ & 1.43 & 1.65 & 1.61 \\
		                                 & $\|u-\uP_{h}\|_{H^{1}}$ & 0.46 & 0.59 & 0.64 \\
		                                 & $\|\lP-\lP_{h}\|_{L^{2}}$ & 1.16 & 1.30 & 2.16 \\
		                                 & $\|\lP-\lP_{h}\|_{H^{1}}$ & 0.46 & 0.30 & 1.24 \\
		\hline
		\multirow{4}{*}{Fully stabilized} & $\|u-\uP_{h}\|_{L^{2}}$ & 1.20 & 1.64 & 1.84 \\
		                                 & $\|u-\uP_{h}\|_{H^{1}}$ & 0.57 & 0.72 & 0.81 \\
		                                 & $\|\lP-\lP_{h}\|_{L^{2}}$ & 1.10 & 1.39 & 1.82 \\
		                                 & $\|\lP-\lP_{h}\|_{H^{1}}$ & 0.60 & 0.82 & 0.99 \\
		\hline
		\multirow{4}{*}{Minimally stabilized} & $\|u-\uP_{h}\|_{L^{2}}$ & 1.16 & 1.51 & 1.82 \\
		                                     & $\|u-\uP_{h}\|_{H^{1}}$ & 0.57 & 0.73 & 0.81 \\
		                                     & $\|\lP-\lP_{h}\|_{L^{2}}$ & 1.11 & 1.48 & 1.93 \\
		                                     & $\|\lP-\lP_{h}\|_{H^{1}}$ & 0.60 & 0.78 & 0.98 \\
		\bottomrule
	\end{tabular}
	\label{tab:scalar_errors}
\end{table}
When comparing the same fields under the equal-order formulation, the convergence rates remain similar in magnitude to those obtained with the Natural formulation, with no substantial improvement. Among the equal-order variants, we highlight the enhanced convergence of $ \uP $ and $ \lP $ in the $ H^1 $-norm achieved by the Minimally and Fully stabilized formulations, compared to the Unstabilized case. However, no further improvements are observed in the $ L^2 $-norm for either field when stabilization is applied.
For the vector fields $ \sP$, $ \eP $, and $ \mP $, we observe that the convergence rates in the $ L^2 $-norm remain nearly identical across all formulations and geometries. Improvements in convergence are primarily associated with increasing the interior angle at the re-entrant corner, with higher accuracy achieved when transitioning from acute to obtuse geometries. 
These results confirm that the convergence behavior is primarily dictated by the corner-induced singularity. While stabilization enhances robustness -- especially in the $H^1$-norm for scalar fields -- it does not overcome the inherent regularity limitations imposed by the geometry. 
\begin{table}[h!]
	\centering
	\caption{Estimated convergence rates in the $ L^2 $-norm for the vector fields $ \sP_h $, $ \eP_h $, and $ \mP_h $ using various stabilization schemes and different reentrant corner angles $ \psi \in \left\{ \frac{7\pi}{4}, \frac{3\pi}{2}, \frac{5\pi}{4} \right\} $.}
	%\rowcolors{3}{black!10!white}{white}
	\begin{tabular}{l l c c c}
		\toprule
		%\rowcolor{black!30!white} 
		Formulation & Norms & $\psi = \frac{7\pi}{4}$ & $\psi = \frac{3\pi}{2}$ & $\psi = \frac{5\pi}{4}$ \\ [2pt]
		\hline
		\multirow{3}{*}{Natural}    & $\|\mathbf{s}-\sP_{h}\|_{L^{2}}$ & 0.57 & 0.73 & 0.83  \\
		                                 & $\|\mathbf{e}-\eP_{h}\|_{L^{2}}$ & 0.57 & 0.72 & 0.81 \\
		                                 & $\|\mP-\mP_{h}\|_{L^{2}}$ & 0.58 & 0.73 & 0.87 \\
		\hline
		\multirow{3}{*}{Unstabilized}    & $\|\mathbf{s}-\sP_{h}\|_{L^{2}}$ & 0.55 & 0.67 & 0.86 \\
		                                 & $\|\mathbf{e}-\eP_{h}\|_{L^{2}}$ & 0.56 & 0.72 & 0.84 \\
		                                 & $\|\mP-\mP_{h}\|_{L^{2}}$ & 0.64 & 0.69 & 0.97 \\
		\hline
		\multirow{3}{*}{Fully stabilized} & $\|\mathbf{s}-\sP_{h}\|_{L^{2}}$ & 0.56 & 0.72 & 0.84 \\
		                                 & $\|\mathbf{e}-\eP_{h}\|_{L^{2}}$ & 0.55 & 0.73 & 0.82 \\
		                                 & $\|\mP-\mP_{h}\|_{L^{2}}$ & 0.55 & 0.73 & 0.89 \\
		\hline
		\multirow{3}{*}{Minimally stabilized} & $\|\mathbf{s}-\sP_{h}\|_{L^{2}}$ & 0.55 & 0.73 & 0.82 \\
		                                     & $\|\mathbf{e}-\eP_{h}\|_{L^{2}}$ & 0.55 & 0.72 & 0.82 \\
		                                     & $\|\mP-\mP_{h}\|_{L^{2}}$ & 0.58 & 0.69 & 0.83 \\
		\bottomrule
	\end{tabular}
	\label{tab:vector_errors}
\end{table}


\newpage
\subsection{Test case 3: data assignation}

In the previous numerical examples, the exact values of $ (\etP, \stP) $ were available at every integration point in $\Omega$, such that we could assess the convergence
properties of the finite element formulation. 
Now, as a final numerical experiment, we assess the consistency of the discrete 
formulation with respect to the number of available data pairs $ (\etP, \stP) $, aiming 
to mimic a more realistic scenario for data-driven applications. 

To this end, we consider a structured cartesian grid of the computational domain into $N_{\mbox{\tiny samples}} = N_d \times N_d$ non-overlapping quadrilateral subdomains $ \{\Omega_i\}_{i=1}^{N_{\mbox{\tiny samples}}} $ where the solution will be sampled. For each subdomain $ \Omega_i $, let $\mathbf{x}_{i}^{c} \in \Omega$ denotes its centroid, at which a unique data pair $(\etP_i, \stP_i)$ is associated. For this problem we consider the manufactured solution: 
\begin{equation}
	\uP: \Omega \rightarrow \mathbb{R}, \quad \uP(x,y) = \sin(\pi x)\sin(\pi y),
\end{equation}
whose gradient is given by:
\begin{equation}
	\eP: \Omega \rightarrow \mathbb{R}^{2}, \quad \eP(x,y) = \pi \, \left[\cos(\pi x)\sin(\pi y), \sin(\pi x)\cos(\pi y) \right]^{\top}.
\end{equation}
For the sake of simplicity, we consider $ \sP = -\eP $ and further assume $\lP  = 0$, $ \mP = \left[0,0\right]^{\top}$,
such that, the data fields $(\etP, \stP)$ coincide with $(\eP, \sP)$ in the exact setting. 
Additionally, for the source term in Eq.~(\ref{eq:dif4}) we have: 
$$
q: \Omega \rightarrow \mathbb{R}, \quad q(x,y) = \left(1 + 2\pi^2 \right)\sin(\pi x)\sin(\pi y).
$$

The values of $\{(\etP_i, \stP_i)\}_{i=1}^{N_{\mbox{\tiny samples}}}$ are thus obtained by evaluating the manufactured fields $(\eP, \sP)$ at $\mathbf{x}_{i}^{c}$. To assign a data pair $ (\etP_i, \stP_i) $ from the dataset to each element of the mesh, we use a geometric proximity criterion, i.e., denoting by $ \mathbf{x}_K^{c} $ the centroid of mesh 
element $ K \in \mathcal{T}_h $, the data point assigned to element $K$ is given by
\begin{equation}
	(\etP_h, \stP_h)(\mathbf{x}_{K}^{c}) := (\etP_i, \stP_i) \quad \text{where} \quad i = \arg\min_{j} \| \mathbf{x}_K^{c} - \mathbf{x}_{j}^{c} \|_2
\end{equation}
where $\| \cdot \|_2$ stands for the Euclidean norm. Notice that, in this way, the data fields $(\etP, \stP)$
have an elementwise constant representation over $\Omega$ in which all the elements of
$\mathcal{T}_h$ that are close to a sampled point share the same value. 
This does not introduce any additional
approximation error only when the vector fields are in $[\DG_0]^2$.
Finally, in the numerical experiments, we consider only $ k=0 $ for both formulations 
and assume pure Dirichlet boundary conditions.

To illustrate the solution obtained with this approach, figure \ref{fig:compound_Nsamples} shows
contour plots of pressure and flux for the Natural formulation 
varying the number of subdomains from $2\times 2$ samples up to $20 \times 20$. 
In the last row of figures, we also show contours of flux for the equal order
formulation, the pressure for this case being very similar to the former. 
The results can be compared to the most refined partition, which encompasses $1000 \times 1000$
samples and is indistinguishable from the exact solution. 
The coarsest samplings produce pressure fields that exhibit a pyramidal shape with
as many faces as number of samples considered. The $20 \times 20$ solution 
already shows a high degree of similarity to the exact one.
\begin{figure}[h!]
	\centering
	\caption{Contours of pressure and flux obtained using the Natural formulation, with an increasing number of data samples distributed over the domain. The last row also shows the corresponding flux field computed using the equal-order formulation. The pressure field for this case is omitted, as it is very
	similar to that obtained with the Natural formulation.}
	\fontsize{8pt}{10pt}\selectfont 
	\def\svgwidth{1.0\columnwidth}
	\input{Figures/PDFs/compound_Nsamples.pdf_tex}
	\label{fig:compound_Nsamples}
\end{figure}

Now, we quantitatively assess the convergence of the method under these conditions for
increasingly refined sampling partitions, from $64 \times 64$ up to $2048 \times 2048$.
The domain is discretized with a unstructured mesh of characteristic size $h$ which
is also successively refined, such that there is no conformity between the sampling
subdomains and the mesh elements. Figures~\ref{fig:quadricula_convergence_natural_linear} to~\ref{fig:quadricula_convergence_equal_order_minimally_stabilized_linear} show the mesh convergence behavior for the scalar fields $\uP$ and $\lP$ in the $L^{2}$- and $H^{1}$-norms, as well as for the vector fields $\eP$, $\sP$, and $\mP$ in the $L^{2}$-norm, for different values of $N_{\mbox{\tiny samples}}$. Figure~\ref{fig:quadricula_convergence_natural_linear} corresponds to the Natural formulation, whereas Figures~\ref{fig:quadricula_convergence_equal_order_unstabilized_linear} to~\ref{fig:quadricula_convergence_equal_order_minimally_stabilized_linear} correspond to the equal order formulation variants: Unstabilized, Fully Stabilized, and Minimally Stabilized, respectively.

For the Natural formulation, the mesh convergence orders approach the theoretical error estimates as the number of sample subdomains $N_{\mbox{\tiny samples}}$ increases. A similar trend is observed for the variants of the equal order formulation. In particular, the field $\uP$ exhibits $\mathcal{O}(h^{2}) $ in the $L^{2}$-norm and $ \mathcal{O}(h) $ in the $H^{1}$-norm in all formulations, with only small variations in the error magnitude.

Similarly, nearly linear convergence rates are observed for the vector fields $\sP$, $\eP$, and $\mP$ in all formulations, with the Natural formulation achieving consistent linear convergence for all three fields. Linear convergence is further achieved for the vector fields with $N_{\mbox{\tiny samples}} = 512 \times 512$ in the Natural formulation, while approximately $N_{\mbox{\tiny samples}}  = 2048 \times 2048$ are required to observe the same convergence behavior for the equal order variants. Despite this, the equal order formulations tend to exhibit smaller error magnitudes.

The first point to note is that the number of data samples limits the convergence rate of the method. Specifically, for a fixed number of data points, refining the finite element mesh eventually leads to error stagnation. However, as the number of samples increases, the method begins to exhibit convergence, at least within the range of $h$ values considered.

As a final observation, the field $\lP$ does not achieve the expected convergence rates in $ H^{1} $-norm, except in the Unstabilized equal order variant. However, increasing $N_{\mbox{\tiny samples}}$ systematically reduces the associated errors, indicating an overall improvement in the approximation quality.

These results already show a consistency with respect to mesh refinement and size of the sample set for the different finite element formulations assessed which makes them a suitable candidates for data-driven applications. 




\begin{figure}[h!]
	\centering
	\caption{Mesh convergence results for the scalar fields $\uP$ and $\lP$ in the $L^2$ and $H^1$ norms, and for the vector fields $\eP$, $\sP$, and $\mP$ in the $L^2$ norm, using the Natural
	 formulation across various numbers of subdomains $N_{\mbox{\tiny samples}}$.}
	%\fontsize{11pt}{10pt}\selectfont 
	\def\svgwidth{1.0\columnwidth}
	\input{Figures/PDFs/quadricula_convergence_natural_linear.pdf_tex}
	\label{fig:quadricula_convergence_natural_linear}
\end{figure}


\begin{figure}[h!]
	\centering
	\caption{Mesh convergence results for the scalar fields $\uP$ and $\lP$ in the $L^2$ and $H^1$ norms, and for the vector fields $\eP$, $\sP$, and $\mP$ in the $L^2$ norm, using the  Unstabilized equal-order formulation across various numbers of subdomains $N_{\mbox{\tiny samples}}$.}
	%\fontsize{11pt}{10pt}\selectfont 
	\def\svgwidth{1.0\columnwidth}
	\input{Figures/PDFs/quadricula_convergence_equal_order_unstabilized_linear.pdf_tex}
	\label{fig:quadricula_convergence_equal_order_unstabilized_linear}
\end{figure}

\begin{figure}[h!]
	\centering
	\caption{Mesh convergence results for the scalar fields $\uP$ and $\lP$ in the $L^2$ and $H^1$ norms, and for the vector fields $\eP$, $\sP$, and $\mP$ in the $L^2$ norm, using the Fully Stabilized equal-order formulation across various numbers of subdomains $N_{\mbox{\tiny samples}}$.}
	%\fontsize{11pt}{10pt}\selectfont 
	\def\svgwidth{1.0\columnwidth}
	\input{Figures/PDFs/quadricula_convergence_equal_order_fully_stabilized_linear.pdf_tex}
	\label{fig:quadricula_convergence_equal_order_fully_stabilized_linear}
\end{figure}


\begin{figure}[h!]
	\centering
	\caption{Mesh convergence results for the scalar fields $\uP$ and $\lP$ in the $L^2$ and $H^1$ norms, and for the vector fields $\eP$, $\sP$, and $\mP$ in the $L^2$ norm, using the Minimally Stabilized equal-order formulation across various numbers of subdomains $N_{\mbox{\tiny samples}}$.}
	%\fontsize{11pt}{10pt}\selectfont 
	\def\svgwidth{1.0\columnwidth}
	\input{Figures/PDFs/quadricula_convergence_equal_order_minimally_stabilized_linear.pdf_tex}
	\label{fig:quadricula_convergence_equal_order_minimally_stabilized_linear}
\end{figure}


% ====== DO NOT DELETE THIS ====== %

%\begin{table}[h!]
%	\centering
%	\begin{tabular}{l l l l l}
%		Formulation & Norms & $\gamma = \frac{\pi}{4}$ & $\gamma = \frac{\pi}{2}$ & $\gamma = \frac{3\pi}{4}$ \\
%		\hline
%		\mPltirow{4}{*}{Natural}				 & $\|u-\uP_{h}\|_{L^{2}}$ & 1.135  & 1.438  & 1.646  \\
%		                                 & $\|u-\uP_{h}\|_{H^{1}}$ & 0.5699 & 0.7171  & 0.8139 \\
%		                                 & $\|\lP-\lP_{h}\|_{L^{2}}$ & 1.217  & 1.684 & 1.965  \\
%		                                 & $\|\lP-\lP_{h}\|_{H^{1}}$ & 0.9898  & 0.9941 & 0.9913  \\
%		\hline
%		\mPltirow{4}{*}{Unstabilized}    & $\|u-\uP_{h}\|_{L^{2}}$ & 1.427 & 1.654 & 1.613 \\
%		                                 & $\|u-\uP_{h}\|_{H^{1}}$ & 0.4591 & 0.5946 & 0.6426 \\
%		                                 & $\|\lP-\lP_{h}\|_{L^{2}}$ & 1.161 & 1.301 & 2.161 \\
%		                                 & $\|\lP-\lP_{h}\|_{H^{1}}$ & 0.4593 & 0.296 & 1.243 \\
%		\hline
%		\mPltirow{4}{*}{Fully stabilized} & $\|u-\uP_{h}\|_{L^{2}}$ & 1.204 & 1.635 & 1.843 \\
%		                                 & $\|u-\uP_{h}\|_{H^{1}}$ & 0.5702 & 0.7178 & 0.8148 \\
%		                                 & $\|\lP-\lP_{h}\|_{L^{2}}$ & 1.102 & 1.39 & 1.817 \\
%		                                 & $\|\lP-\lP_{h}\|_{H^{1}}$ & 0.6008 & 0.824 & 0.9891 \\
%		\hline
%		\mPltirow{4}{*}{Minimally stabilized} & $\|u-\uP_{h}\|_{L^{2}}$ & 1.155 & 1.514 & 1.817 \\
%		                                     & $\|u-\uP_{h}\|_{H^{1}}$ & 0.567 & 0.7261 & 0.8138 \\
%		                                     & $\|\lP-\lP_{h}\|_{L^{2}}$ & 1.108 & 1.479 & 1.93 \\
%		                                     & $\|\lP-\lP_{h}\|_{H^{1}}$ & 0.597 & 0.7779 & 0.9834 \\
%		\hline
%	\end{tabular}
%	\label{tab:condkjfs}
%	\caption{Convergence rates for the scalar variables $u$ and $\lP$ under different stabilization schemes.}
%\end{table}

%\begin{table}[h!]
%	\centering
%	\begin{tabular}{l l c c c}
%		Formulation & Norms & $\gamma = \frac{\pi}{4}$ & $\gamma = \frac{\pi}{2}$ & $\gamma = \frac{3\pi}{4}$ \\
%		\hline
%		\mPltirow{3}{*}{Natural}    & $\|\mathbf{s}-\sP_{h}\|_{L^{2}}$ & 0.5676 & 0.7252 & 0.8269  \\
%		                                 & $\|\mathbf{e}-\eP_{h}\|_{L^{2}}$ & 0.5698 & 0.7171 & 0.8138 \\
%		                                 & $\|\mP-\mP_{h}\|_{L^{2}}$ & 0.5817 & 0.7298 & 0.8652 \\
%		\hline
%		\mPltirow{3}{*}{Unstabilized}    & $\|\mathbf{s}-\sP_{h}\|_{L^{2}}$ & 0.5529 & 0.6674 & 0.8582 \\
%		                                 & $\|\mathbf{e}-\eP_{h}\|_{L^{2}}$ & 0.5556 & 0.7233 & 0.8379 \\
%		                                 & $\|\mP-\mP_{h}\|_{L^{2}}$ & 0.6398 & 0.6906 & 0.9706 \\
%		\hline
%		\mPltirow{3}{*}{Fully stabilized} & $\|\mathbf{s}-\sP_{h}\|_{L^{2}}$ & 0.5601 & 0.7196 & 0.8416 \\
%		                                 & $\|\mathbf{e}-\eP_{h}\|_{L^{2}}$ & 0.5495 & 0.7262 & 0.821 \\
%		                                 & $\|\mP-\mP_{h}\|_{L^{2}}$ & 0.5543 & 0.7329 & 0.8948 \\
%		\hline
%		\mPltirow{3}{*}{Minimally stabilized} & $\|\mathbf{s}-\sP_{h}\|_{L^{2}}$ & 0.5496 & 0.7266 & 0.8239 \\
%		                                     & $\|\mathbf{e}-\eP_{h}\|_{L^{2}}$ & 0.5537 & 0.7209 & 0.8224 \\
%		                                     & $\|\mP-\mP_{h}\|_{L^{2}}$ & 0.5820 & 0.6921 & 0.8324 \\
%		\hline
%	\end{tabular}
%	\label{tab:vector_errors}
%	\caption{Convergence rates for the vector variables $\mathbf{s}$, $\mathbf{e}$, and $\mP$ under different stabilization schemes.}
%\end{table}


